\documentclass[twoside]{article}
\usepackage[utf8]{inputenc}
\usepackage{listings}
\usepackage[dvipsnames]{xcolor}


% The package supports other formats as well, but this template
% is primarily checked and updated for examinations at TUM BGDM

% with/without answers
% \usepackage[uebung, answers]{tumbgdm}
\usepackage[uebung]{tumbgdm}
% \setlength{\parskip}{\baselineskip}%
% \setlength{\parindent}{0pt}%
% \usepackage{pythonhighlight}
\usepackage[skip=6pt, indent=0pt]{parskip}
\usepackage{xcolor}
\hypersetup{
    colorlinks,
    linkcolor={red!50!black},
    citecolor={blue!50!black}, 
    urlcolor={blue!80!black}
}

\Professor{Prof. Dr. Martin Werner}
\Universitaet{Technical University of Munich}
\Institut{School of Engineering and Design}
\Lehrstuhl{Professorship Big Geospatial Data Management}
\Semester[WiSe 26/27]{WiSe 26/27}
\Vorlesung{Principles of Programming}
\InstitutsWebsite{https://www.bgd.ed.tum.de/}

\nummerUebungsblatt{1}
\nummerErsteFrage{1}


\LearningOutcome{Basic Computing}


\begin{document}
\maketitle

\section{Python Loops}

\begin{task}{For Loops}{}{}
Use one or more \texttt{for} loops to generate and print the multiplication table for the numbers 1 through 10. 
\begin{itemize} 
    \item The rows should represent the multiplicands (1--10). 
    \item The columns should represent the multipliers (1--10). 
    \item Each cell should contain the product of the row and column values. 
    \item Do not manually print each row. Generate the table using loops. 
\end{itemize} \vspace{10pt}

\vspace{10pt}
Expected output:
\begin{lstlisting}[numbers=none]
1   2  3  4  5  6  7  8  9  10 		
2   4  6  8 10 12 14 16 18  20 		
3   6  9 12 15 18 21 24 27  30 		
4   8 12 16 20 24 28 32 36  40 		
5  10 15 20 25 30 35 40 45  50 		
6  12 18 24 30 36 42 48 54  60 		
7  14 21 28 35 42 49 56 63  70 		
8  16 24 32 40 48 56 64 72  80 		
9  18 27 36 45 54 63 72 81  90 		
10 20 30 40 50 60 70 80 90 100 
\end{lstlisting}

\end{task}

\newpage
\begin{task}{While Loops}{}{}
Write a function that takes two integer inputs: 
\begin{itemize} 
    \item \texttt{base} 
    \item \texttt{exponent} 
\end{itemize} 

Use a \texttt{while} loop to calculate the value of $ base^{exponent}$ using repeated multiplication.

\textbf{Do not use} Python's built-in \texttt{pow()} function or the \texttt{**} operator.

Print the result in the format shown below.

\vspace{10pt}
\textbf{Example}\\
Expected output for input base=2 and exp=5:
\begin{lstlisting}[numbers=none]
2 to the power of 5 is:  32
\end{lstlisting}

\end{task}

\newpage
\section{Python List Comprehensions}

% Solve the following tasks with list comprehensions:
% \begin{lstlisting}[numbers=none]
% newlist = [expression for item in iterable if condition] 
% \end{lstlisting}


\begin{task}{Reverse String}{}{}
Using a list comprehension, create a function that returns the characters of a string in reverse order.
The function should return a list containing the reversed characters.

\vspace{10pt}
\textbf{Example}\\
Expected output for input s = "hello":
\begin{lstlisting}[numbers=none]
["o","l","l","e","h"]
\end{lstlisting}

\textbf{Note:} Use a list comprehension rather than Python's built-in \texttt{reverse()} method or slicing syntax

%\vspace{5pt}
%\textbf{Constraints:}
%\begin{itemize}
%    \item $1 \leq len(s) \leq 10^5$
%    \item s[i] is a printable ascii character.
%\end{itemize}
\end{task}


\begin{task}{Sort Array}{}{}
Given an integer array \texttt{nums}, create a new array in which all even numbers appear first, followed by all odd numbers. 
Use one or more list comprehensions to construct the result. 
The relative order of the even numbers and odd numbers does not matter.

\vspace{10pt}
\textbf{Example}\\
Expected output for input nums = [3,1,2,4]:
\begin{lstlisting}[numbers=none]
[2,4,3,1]
\end{lstlisting}
Explanation: The outputs [4,2,3,1], [2,4,1,3], and [4,2,1,3] would also be accepted.

%\textbf{Constraints:}
%\begin{itemize}
%    \item $1 \leq len(nums) \leq 5000$
%    \item $0 \leq nums[i] \leq 5000$
%\end{itemize}
\end{task}


\begin{task}{Fizz Buzz}{}{}
Fizz buzz is a group word game for children to teach them about division. Players take turns to count incrementally, replacing any number divisible by three with the word "fizz", and any number divisible by five with the word "buzz", and any number divisible by both 3 and 5 with the word "fizzbuzz" (from \href{https://en.wikipedia.org/wiki/Fizz_buzz}{Wikipedia}).\\

Given an integer n, return a string array answer (1-indexed) where:

\begin{enumerate}
    \item answer[i] == “FizzBuzz” if i is divisible by 3 and 5.
    \item answer[i] == “Fizz” if i is divisible by 3.
    \item answer[i] == “Buzz” if i is divisible by 5.
    \item answer[i] == i (as a string) if none of the above conditions are true.
\end{enumerate}

\vspace{10pt}
\textbf{Example}\\
Expected output for input n = 15:
\begin{lstlisting}[numbers=none]
["1","2","Fizz","4","Buzz","Fizz","7","8","Fizz","Buzz","11","Fizz","13","14","FizzBuzz"]
\end{lstlisting}

%\vspace{5pt}
%\textbf{Constraints:}
%\begin{itemize}
%    \item $1 \leq n \leq 10^4$
%\end{itemize}
\end{task}

\begin{task}{Numpy}{}{}
Import the NumPy package (official documentation: \url{https://numpy.org/doc/}) and use it to solve the following tasks:

\begin{enumerate} 
    \item Print the NumPy version and configuration information. 
    \item Create a 10×10 matrix filled with zeros. 
    \item Determine and print the shape of an array (e.g., the matrix created in Task 1). 
    \item Create a 4×6 matrix filled with zeros and set the element at index \texttt{[2, 3]} to 1. 
    \item Create a 3×7 matrix containing consecutive integers starting from 21. 
    \item Create a vector of length 10 filled with random values between 0 and 1. 
    \item Reverse the vector from Task 6 (i.e., the first element becomes the last). 
    \item Create a 3×3 matrix containing random integer values between 0 and 7. 
    \item Find the indices of all non-zero elements in: 
        \begin{lstlisting}[numbers=none] 
[1, 0, 7, 0, 1, 0, 0, 1, 0, 0, 1] 
        \end{lstlisting} 
    \item Find all unique values in: 
    \begin{lstlisting}[numbers=none] 
[1, 0, 7, 0, 1, 0, 0, 1, 0, 0, 1] 
    \end{lstlisting} 
    \item Create a 3×3 identity matrix. 
    \item Multiply a 5×3 matrix by a 3×2 matrix. 
    \item Create a random vector of length 10 and sort it in ascending order. 
    \item Create a 5×5×5 array of random values and determine its minimum, maximum, and mean values. 

\end{enumerate}
\end{task}

\end{document}
